\section{The Ideal Problem}
\label{sec:ideal}

This section defines the compilation problem: a contract over observable
behavior, the programs that satisfy it, and an objective that picks a cheapest
one. Leroy's preservation of a
specification~\cite[\S2.1]{leroy2009backend} and the refinement
calculi~\cite{back1998refinement,morgan1994programming} are the setting. Named
implementation parameters, a taxonomy of slack, that objective, and then a
split between the functional contract and the optimization contract make the
problem specific enough to compile against.

The definitions apply to any language with a deterministic sequential semantics
and a notion of observable behavior. The running examples are in C, and the
as-if rule and undefined behavior are the same in C++.

\subsection{Programs and the As-If Rule}
\label{sec:setting}

Let \(\Sigma\) be the set of initial states of a program, covering its inputs
and initial memory. Running program \(P\) on \(\sigma \in \Sigma\) has one of
three results. It may terminate in a final state \(\sigma'\) after producing a
trace \(\tau \in \mathcal{T}\) of observable events, written
\(\langle P, \sigma \rangle \Downarrow (\sigma', \tau)\). It may run forever,
written \(\mathsf{div}\). Or it may reach an operation whose meaning the
language leaves undefined, written \(\mathsf{ub}\). We write
\(\llbracket P \rrbracket(\sigma)\) for the result.

In C and C++, the observable behavior consists of I/O and accesses to volatile
objects~\cite[\S4.1.2]{n4950}. Local values, the order of intermediate
computations, and unescaped layout are not observable, so comparisons use only
those observable parts.

\begin{definition}[Observational Equivalence]
\(P_1 \approx P_2\) iff \(\llbracket P_1 \rrbracket(\sigma)\) and
\(\llbracket P_2 \rrbracket(\sigma)\) have the same observable part for every
\(\sigma \in \Sigma\). We write \([P]_\approx\) for the set of programs
observationally equivalent to \(P\).
\end{definition}

\([P]_\approx\) is the space a conventional optimizing compiler may search.
Undefined behavior loosens this slightly: on inputs where the source is
undefined, the compiled program may do anything~\cite{lopes2021alive2,
lee2017taming}.

\subsection{Named Parameters and Functional Contracts}
\label{sec:contract}

Implementation parameters are named constants in the source, written
\(\theta_P\). The set of names is part of the contract. Extracting
\(\theta_P\) is syntactic: a \texttt{\#define} capacity, not a property of
\(\llbracket P \rrbracket\). Two observationally equivalent programs may
differ in \(\theta\). The contract is indexed by \(\theta\), so
\(\mathit{Post}_\theta\) can depend on the capacity the implementation chose.

\begin{definition}[Functional Intent Contract]
\label{def:contract}
A functional intent contract is a triple
\[
\Phi_{\text{func}} = \langle \mathit{Pre}, \mathit{Post}, \mathcal{I} \rangle
\]
where \(\mathit{Pre} \subseteq \Sigma\) is the set of initial states the
program must handle, \(\mathcal{I}\) is the set of acceptable parameter
valuations, and \(\mathit{Post}\) gives, for each \(\theta \in \mathcal{I}\),
a relation \(\mathit{Post}_\theta\) between an initial state and the
acceptable observable results.
Program \(P\) satisfies \(\Phi_{\text{func}}\), written
\(P \models \Phi_{\text{func}}\), iff \(\theta_P \in \mathcal{I}\) and
\[
\forall \sigma \in \mathit{Pre}.\ \exists \sigma', \tau.\
\langle P, \sigma \rangle \Downarrow (\sigma', \tau)
\ \wedge\ (\sigma, (\sigma', \tau)) \in \mathit{Post}_{\theta_P}.
\]
\end{definition}

The triple follows Design by Contract~\cite{meyer1988eiffel}.
Satisfaction requires termination without undefined behavior on every input
in \(\mathit{Pre}\), and the contract mentions only observable results. The
set of satisfying programs is
\(\mathrm{Sat}(\Phi_{\text{func}}) = \{P \mid P \models \Phi_{\text{func}}\}\).
We write \(P_1 \sim_{\Phi_{\text{func}}} P_2\) when both lie in that set.

Each component states one kind of intent. \(\mathit{Pre}\) states which
inputs matter. \(\mathcal{I}\) states acceptable parameter valuations, so a
queue that must absorb bursts of 800 packets with 25\% headroom has
\(\mathcal{I} = \{\theta \mid \theta_{\mathrm{cap}} \ge 1000\}\).
\(\mathit{Post}_\theta\) states what the program must do, relative to the
parameters it chose: enqueued packets come out in order, and an enqueue fails
exactly when \(\theta_{\mathrm{cap}}\) packets are stored.

\begin{definition}[Realizability]
\(\Phi_{\text{func}}\) is vacuous if \(\mathit{Pre} = \emptyset\). It is
contradictory if no parameter valuation, input, and result jointly satisfy
\(\mathit{Post}\). It is realizable if there exists \(\theta \in \mathcal{I}\)
such that every \(\sigma \in \mathit{Pre}\) has at least one acceptable
result under \(\mathit{Post}_\theta\).
\end{definition}

Vacuous contracts are rejected before search: they have empty
\(\mathit{Pre}\). A contract whose \(\mathit{Post}\) no program can satisfy
is discovered later, when no candidate is proved. Realizability is the
\(\forall\sigma\,\exists o\) condition.

\subsection{Slack}
\label{sec:slack}

\(\mathrm{Sat}(\Phi_{\text{func}})\) differs from \([P]_\approx\) where the
contract leaves a decision open. There are four ways to leave one open.

\begin{itemize}
  \item \emph{Input slack.} \(\mathit{Pre} \neq \Sigma\). On excluded inputs
  any behavior is acceptable. This is the kind undefined behavior already
  provides, and the kind SeGaBench grants by fixing a valid-input
  domain~\cite{jiang2026segabench}.
  \item \emph{Output slack.} For some covered input, \(\mathit{Post}_\theta\)
  accepts two distinct observable results. A process-ID allocator must return
  an unused ID, or \(\texttt{NONE}\) when every ID is in use. Scanning from
  ID 1, and continuing from a cursor, both satisfy that, and they return
  different IDs once an ID below the cursor has been freed. No conventional
  compiler may replace one by the other.
  \item \emph{Logical-parameter slack.} The contract accepts several
  parameter valuations, and the choice is visible in the results. Queues of
  capacity 1000 and 1024 reject a full-queue enqueue at different points. Both
  are acceptable when \(\mathcal{I}\) requires only a capacity of at least
  1000.
  \item \emph{Unexposed slack.} Extra resources or a different representation
  that the interface does not expose. A queue that allocates 1024 slots and
  still becomes full at 1000 uses a power-of-two mask while preserving
  caller-facing capacity. Folly and DPDK's exact-size rings are this
  case~\cite{folly_pcq,dpdk_rte_ring}. If the extra cells and the mask are
  unobservable at the API, the as-if rule may already permit the change,
  provided the analysis can see it. The contract still has to license the
  extra memory, which is a budget concern
  (\S\ref{sec:opt_contract}).
\end{itemize}

Input and output slack are don't-cares in Karnaugh's sense: an incompletely
specified function can be minimized by assigning combinations the
specification leaves open~\cite{karnaugh1953map}.
Logical-parameter slack and unexposed slack concern the implementation
rather than unspecified I/O. The first is a named parameter visible in the
results, and the second is extra store or a representation the interface
does not name.

The four kinds of slack are properties of \(\Phi_{\text{func}}\). Whether
\(\mathrm{Sat}(\Phi_{\text{func}})\) is larger than \([P]_\approx\) is
settled by local patching.

A language supports \emph{local patching} if for every program \(P\), state
\(\sigma_0\), and producible observable result \(o\), some program behaves
like \(P\) everywhere except at \(\sigma_0\), where it produces \(o\). Any
language that can test its input against a constant has this property,
including essentially every general-purpose language in use today.

\begin{proposition}[Legal Set Expands Exactly on Observable Slack]
\label{prop:expansion}
Assume local patching, and that each \(\mathit{Post}_\theta\) with
\(\theta \in \mathcal{I}\) is realized by some program. Let
\(P \models \Phi_{\text{func}}\). Then
\(\mathrm{Sat}(\Phi_{\text{func}})\) contains a program that is not
observationally equivalent to \(P\) iff the contract has input, output, or
logical-parameter slack, counting only results some program produces.
Unexposed slack is not in that list: it need not change those results.
\end{proposition}

\noindent\emph{Proof sketch.} If the contract accepts a result on some input
that differs from \(P\)'s, patching produces a legal program outside
\([P]_\approx\). Without slack, every covered input has one acceptable result,
the same for every allowed parameter, and every input is covered, so every
legal program matches \(P\). (Appendix~\ref{app:proofs}.)

When the contract has no input, output, or logical-parameter slack, intent
compilation is conventional compilation: the approach generalizes the as-if
rule rather than replacing it. Unexposed slack can still change cost and
memory. The proposition does not say that any extra legal program is faster.

\subsection{The Ideal Objective}
\label{sec:ideal_obj}

Let \(\mathcal{L}_{\text{Target}}\) be an unmodified production compiler for
the target. It implements the as-if rule: idealized, it searches
\([P]_\approx\). Let \(\mathcal{O}(\mathrm{Bin}, \text{Target})\) rank
binaries by cost on that target. A conventional compiler, idealized,
returns a cheapest binary in the as-if class of \(P_1\):
\begin{equation}
\label{eq:traditional}
\mathrm{Bin}^*_{\text{trad}} \in
\arg\min_{\mathrm{Bin} \in \{\mathcal{L}_{\text{Target}}(P') \mid
P' \in [P_1]_\approx\}}
\mathcal{O}(\mathrm{Bin}, \text{Target}).
\end{equation}
An intent compiler wants the same ranking over every binary that
satisfies the contract:
\begin{equation}
\label{eq:intent_pure}
\mathrm{Bin}^*_{\text{intent}} \in
\arg\min_{\mathrm{Bin} \models \Phi_{\text{func}}}
\mathcal{O}(\mathrm{Bin}, \text{Target}).
\end{equation}
No compiler searches either set as a whole. Conventional compilation
applies heuristic rewrites inside \([P_1]_\approx\), the passes of
\(\mathcal{L}_{\text{Target}}\), and may never reach
\(\mathrm{Bin}^*_{\text{trad}}\). Intent compilation has the same gap on a
larger set, and splits~\eqref{eq:intent_pure} into a source step and an
as-if step:
\begin{align}
\label{eq:intent_source}
P^\star &\in \arg\min_{P' \models \Phi_{\text{func}}}
\mathcal{O}(\mathcal{L}_{\text{Target}}(P'), \text{Target}), \\
\label{eq:intent_lower}
\mathrm{Bin}^*_{\text{intent}} &\in
\arg\min_{\mathrm{Bin} \in \{\mathcal{L}_{\text{Target}}(P'') \mid
P'' \in [P^\star]_\approx\}}
\mathcal{O}(\mathrm{Bin}, \text{Target}).
\end{align}
The first level rewrites the source toward a cheaper program in
\(\mathrm{Sat}(\Phi_{\text{func}})\), and the second is~\eqref{eq:traditional}
with \(P^\star\) in place of \(P_1\). That composition does not guarantee a
minimum of~\eqref{eq:intent_pure}, but when \(P_1\) satisfies the contract it
is never worse than as-if compilation of \(P_1\), because the first level
searches a set that contains every as-if variant of \(P_1\). It separates
program-to-program editing under the contract from semantic lowering of a
single source, and the second level still yields a contract-satisfying
binary from a contract-satisfying source.

\begin{proposition}[As-If Programs Remain Legal]
\label{prop:refine}
If \(P \models \Phi_{\text{func}}\), \(P' \approx P\), and
\(\theta_{P'} = \theta_P\), then \(P' \models \Phi_{\text{func}}\).
\end{proposition}

\noindent\emph{Proof sketch.} A violation by \(P'\) would be a bad parameter
or a bad observable result on some admissible input. The parameters match, so
the first case is impossible. Observational equivalence transfers the second
case to \(P\), contradicting \(P \models \Phi_{\text{func}}\).
(Appendix~\ref{app:proofs}.)

The side condition \(\theta_{P'} = \theta_P\) is required because
\(\mathit{Post}\) is indexed by \(\theta\). Hence the second level
preserves the contract: if \(P^\star \models \Phi_{\text{func}}\), every
program in \([P^\star]_\approx\) with the same \(\theta\) does too.

\subsection{Optimization Contracts}
\label{sec:opt_contract}

The source step~\eqref{eq:intent_source} and the as-if
step~\eqref{eq:intent_lower} are not ranked by the same objective, even
though both displays used \(\mathcal{O}\). The as-if step is
\(\mathcal{L}_{\text{Target}}\)'s search inside \([P^\star]_\approx\).
The source step compares programs that need not be observationally
equivalent, which that backend does not rank. A ranking
for~\eqref{eq:intent_source} has to be supplied as part of the intent.

That ranking searches a larger set than \([P_1]_\approx\). A rewrite in
that set can use more resources than \(P_1\), exceeding what the developer
allocated for the program, for instance a larger memory footprint. Those
allocations are implicit in \(P_1\), so as-if cannot raise them and
\(\mathrm{Sat}(\Phi_{\text{func}})\) does not inherit them. Therefore, an
operational budget \(\mathcal{B}\) is also needed, stating which
implementations are valid to search: those whose resource use stays within
what the developer is willing to allocate to the program.

\begin{definition}[Legality and Feasibility]
\label{def:legal}
A program \(P\) is \emph{legal} iff \(P \models \Phi_{\text{func}}\), and
\emph{feasible} iff \(P \models \mathcal{B}\).
\end{definition}

Search is over programs that are legal and feasible.
\(\mathcal{B}\) and \(\mathcal{O}\) stay outside \(\Phi_{\text{func}}\)
because legality, feasibility, and ranking are different jobs.
\(\Phi_{\text{func}}\) states which programs are legal, \(\mathcal{B}\)
which of those are feasible, and \(\mathcal{O}\) ranks them. That ranking
is supplied as part of the intent. It may be computed from the program
text or measured on a profiled run. \textsc{IntC} uses wall-clock time on
a named driver. Cyclomatic complexity occupies the same slot
(\S\ref{sec:related}). A missed latency target is then suboptimal rather
than illegal, and folding the target into the functional contract would
reject a legal baseline. Budgets describe the machine and change with the
deployment, while \(\Phi_{\text{func}}\) does not.
One functional contract can then yield a 1024-slot ring on a server and a
1000-slot ring under a tight memory budget. A check of \(\mathcal{B}\) on source
approximates a constraint on the compiled binary: padding, alignment, and
allocator overhead are decided later.

Adding \(\mathcal{B}\) as a constraint on the source
step~\eqref{eq:intent_source} gives
\begin{equation}
\label{eq:intent}
P^\star \in \arg\min_{\substack{P' \models \Phi_{\text{func}} \\ P' \models \mathcal{B}}}
\mathcal{O}(\mathcal{L}_{\text{Target}}(P'), \text{Target}),
\qquad
\mathrm{Bin}^*_{\text{intent}} = \mathcal{L}_{\text{Target}}(P^\star),
\end{equation}
where the second equality is~\eqref{eq:intent_lower} once
\(\mathcal{L}_{\text{Target}}\) has searched \([P^\star]_\approx\), and
Section~\ref{sec:compiler} searches a generated subset of that source set.

\begin{definition}[Intent Specification]
\label{def:intent_spec}
An intent specification is a triple
\[
\Phi = \langle \Phi_{\text{func}}, \mathcal{B}, \mathcal{O} \rangle
\]
where \(\Phi_{\text{func}}\) is a functional intent contract
(Definition~\ref{def:contract}), \(\mathcal{B}\) is a conjunction of
operational constraints on a candidate (static memory, a ban on dynamic
allocation, alignment), and \(\mathcal{O}\) is a cost ranking of binaries on
a target.
\end{definition}

If \(P_1\) already satisfies the contract and the budget, the minimum
in~\eqref{eq:intent} is never worse than the one in~\eqref{eq:traditional}
among programs that keep \(\theta_{P_1}\), by Proposition~\ref{prop:refine}.
